% TEMPLATE-MILC-v3.tex - plantilla maestra UNICA de las guias MILC v3.
%
% La usan las tres familias de guias, cada una con su driver:
%   · Grados 6-11  -> scripts/build-guias-g11.py
%   · Bebras       -> scripts/build-guias-bebras.py
%   · Semillero    -> scripts/build-guias-semillero.py
%
% Antes habia tres copias casi identicas (~400 lineas iguales, 6 distintas):
% cada mejora habia que aplicarla tres veces, y en la practica no se hacia
% (el motor de recursos vivia solo en dos de ellas). Lo que varia entre
% familias esta parametrizado en 6 placeholders que cada driver llena
% (se nombran aqui SIN su sintaxis real: el builder reemplaza por texto plano,
% tambien dentro de los comentarios, y un valor multilinea rompe el preambulo):
%
%   HEADER_IZQ        encabezado izquierdo de pagina
%   HEADER_DER        encabezado derecho de pagina
%   PORTADA_ETIQUETA  rotulo sobre el numero grande (GRADO/MOMENTO/MODULO)
%   PORTADA_SERIE     linea de serie de la portada
%   PORTADA_ID        identificador de la guia en la portada
%   PORTADA_CIRCULO   el circulito de grado (°), o vacio si no aplica
%
% El resto de claves las documenta content/guias/_SCHEMA.md. El script valida
% que no queden placeholders sin reemplazar antes de escribir el archivo final.

\documentclass[11pt,letterpaper]{article}
\usepackage[margin=1.75cm]{geometry}
\usepackage[table]{xcolor}
\usepackage{fontspec}
\usepackage[spanish]{babel}
\usepackage{tikz}
% Librerías TikZ para los diagramas de las guías (bloque `recursos:`):
%  · babel  → babel en español vuelve activos caracteres como «>», y TikZ los usa
%             en sus opciones (>=stealth, ->). Sin esta librería, cualquier
%             diagrama con flechas revienta con «\language@active@arg> has an
%             extra }». Debe ir ANTES de que se dibuje nada.
%  · positioning        → sintaxis «below left=2pt and 3pt».
%  · decorations.pathreplacing → llaves ({brace}) para señalar rangos.
\usetikzlibrary{babel,positioning,decorations.pathreplacing}
\usepackage{graphicx}
% Circuitos: el trabajo de electrónica se hace hoy con micro:bit y sus sensores
% integrados (MakeCode, sin cableado), así que no se necesitan esquemáticos.
% Si algún día entran montajes con protoboard, basta descomentar esta línea:
% los diagramas .tex/.tikz declarados en `recursos:` ya se incrustan solos.
% \usepackage{circuitikz}
\usepackage[most]{tcolorbox}
\usepackage{tabularx}
\usepackage{array}
\usepackage{fancyhdr}
\usepackage{titlesec}
\usepackage{ragged2e}
\usepackage{enumitem}
\usepackage{microtype}

% Helvetica Neue no trae la flecha U+2192, asi que cada «→» escrito literalmente
% en el YAML se compilaba a NADA, en silencio (462 flechas en 61 guias: rutas del
% tipo «paleta → tipografia → lockup» salian rotas). Mapeamos el caracter a su
% version matematica, que si tiene glifo. Asi el autor puede seguir escribiendo
% «→» comodamente en el YAML.
\usepackage{newunicodechar}
\newunicodechar{→}{\ensuremath{\rightarrow}}

\setmainfont{Helvetica Neue}
\setsansfont{Helvetica Neue}
\setlength{\parindent}{0pt}
\setlength{\parskip}{6pt}
\hyphenpenalty=200
\exhyphenpenalty=200
\pretolerance=400
\tolerance=2000
\setlength{\emergencystretch}{1em}
\renewcommand{\arraystretch}{1.25}
\setlist{leftmargin=5mm,itemsep=1mm,topsep=1mm}
\newcolumntype{Y}{>{\RaggedRight\arraybackslash}X}

\definecolor{milcMagenta}{HTML}{E6007E}
\definecolor{milcVino}{HTML}{5A0038}
\definecolor{milcTurquesa}{HTML}{0093A5}
\definecolor{milcVerde}{HTML}{58A923}
\definecolor{milcMostaza}{HTML}{E5B400}
\definecolor{milcOcre}{HTML}{B86600}
\definecolor{milcNegro}{HTML}{241816}
\definecolor{milcGris}{HTML}{F7F7F4}
\definecolor{milcLinea}{HTML}{E8E2DD}
\definecolor{milcGradePrimary}{HTML}{84CC16}
\definecolor{milcGradeDark}{HTML}{4D7C0F}
\definecolor{milcGradeSoft}{HTML}{ECFCCB}
\definecolor{milcGradeAccent}{HTML}{CFE7FF}
\definecolor{milcGradeText}{HTML}{FFFFFF}
\color{milcNegro}

\pagestyle{fancy}
\fancyhf{}
\lhead{\small Tecnología e Informática · Grado 7}
\rhead{\small Guía 15 · MILC v3}
\cfoot{\small\thepage}
\renewcommand{\headrulewidth}{0pt}

\newtcolorbox{titlebox}[1]{
  enhanced,colback=#1,colframe=#1,arc=7mm,boxrule=0pt,
  left=7mm,right=7mm,top=3mm,bottom=3mm,
  before skip=8pt,after skip=8pt,
  fontupper=\sffamily\bfseries\Large\color{white}
}

\newtcolorbox{softbox}[3]{
  enhanced,colback=#2,colframe=#1,borderline west={4pt}{0pt}{#1},
  arc=2mm,boxrule=.4pt,left=4mm,right=4mm,top=3mm,bottom=3mm,
  before skip=6pt,after skip=8pt,title={#3},coltitle=white,
  colbacktitle=#1,fonttitle=\sffamily\bfseries,fontupper=\RaggedRight,
  attach boxed title to top left={xshift=3mm,yshift=-2mm},
  boxed title style={arc=2mm, boxrule=0pt}
}

\newtcolorbox{drawbox}[2]{
  enhanced,colback=white,colframe=#1,arc=4mm,boxrule=1pt,
  left=5mm,right=5mm,top=4mm,bottom=4mm,before skip=8pt,after skip=8pt,
  title={#2},coltitle=white,colbacktitle=#1,fonttitle=\sffamily\bfseries,
  fontupper=\RaggedRight,
  attach boxed title to top left={xshift=4mm,yshift=-2mm},
  boxed title style={arc=3mm, boxrule=0pt}
}

\newtcolorbox{infoband}[1]{
  enhanced,colback=#1!8,colframe=#1!50,arc=2mm,boxrule=.5pt,
  left=4mm,right=4mm,top=2mm,bottom=2mm,before skip=4pt,after skip=4pt,
  fontupper=\small\RaggedRight
}

% Puente narrativo entre secciones (contrato editorial Conéctate).
% Da continuidad: "Acabas de X. Ahora vas a Y porque Z."
% Visualmente: banda izquierda en color del grado, texto en itálica pequeña.
\newtcolorbox{puentebox}{
  enhanced,colback=milcGradeSoft,colframe=milcGradeDark,
  borderline west={3pt}{0pt}{milcGradeDark},
  arc=0mm,boxrule=0pt,left=5mm,right=4mm,top=2.5mm,bottom=2.5mm,
  before skip=4pt,after skip=8pt,
  fontupper=\itshape\small\color{milcNegro}\RaggedRight
}
% La flecha va en modo matematico a proposito: Helvetica Neue NO tiene el glifo
% U+2192, asi que \textrightarrow se compilaba a NADA («Missing character: There
% is no → in font Helvetica Neue») y el puente salia sin flecha. En modo
% matematico la flecha viene de la fuente matematica y si se dibuja.
\newcommand{\puente}[1]{%
  \begin{puentebox}%
    \textcolor{milcGradeDark}{$\rightarrow$}\hspace{2mm}#1%
  \end{puentebox}%
}

\newcommand{\linead}[1]{\noindent\color{milcLinea}\rule{#1}{0.5pt}\color{milcNegro}}
\newcommand{\respuesta}[1]{\par\vspace{2mm}\foreach \n in {1,...,#1}{\noindent\color{milcLinea}\rule{\linewidth}{0.5pt}\par\vspace{5mm}}\color{milcNegro}}
\newcommand{\checkbox}{\textcolor{milcGradeDark}{\fbox{\phantom{x}}}\hspace{5pt}}

\newcommand{\phasecard}[4]{
\begin{tcolorbox}[enhanced,colback=#3,colframe=#2,arc=3mm,boxrule=.5pt,height=3.05cm,valign=center,halign=center,left=1.5mm,right=1.5mm,top=2mm,bottom=2mm]
{\sffamily\bfseries\fontsize{22}{24}\selectfont\textcolor{#2}{#1}}\par
{\sffamily\bfseries\small #4}
\end{tcolorbox}
}

% ── Recursos visuales de la guía ──────────────────────────────────────
% Declarados en el bloque `recursos:` del YAML y emitidos por el builder.
% \guiaFigura   → imagen rasterizada o PDF (foto, carta celeste, montaje).
% \guiaDiagrama → fuente TikZ/circuitikz (.tex) incrustada como vector:
%                 es la vía para circuitos y esquemáticos editables.
% #1 = ruta relativa al .tex · #2 = pie (puede ir vacío)
\newcommand{\guiaFigura}[2]{%
\begin{center}
\includegraphics[width=0.88\linewidth,height=0.38\textheight,keepaspectratio]{#1}\\[1.5mm]
{\footnotesize\itshape\textcolor{milcNegro!75}{#2}}
\end{center}
}

\newcommand{\guiaDiagrama}[2]{%
\begin{center}
\input{#1}\\[1.5mm]
{\footnotesize\itshape\textcolor{milcNegro!75}{#2}}
\end{center}
}

\begin{document}
\thispagestyle{empty}
\begin{tikzpicture}[remember picture,overlay]
  \fill[white] (current page.south west) rectangle (current page.north east);
  \path[fill=milcGradePrimary,rounded corners=16mm]
    ([xshift=.55cm,yshift=-.55cm]current page.north west) rectangle
    ([xshift=-.55cm,yshift=3.65cm]current page.south east);

  \node[anchor=north west,text=milcGradeText] at ([xshift=2.0cm,yshift=-1.85cm]current page.north west)
    {\sffamily\bfseries\fontsize{14}{16}\selectfont GRADO};
  \node[anchor=north west,text=milcGradeText,opacity=.82] at ([xshift=2.0cm,yshift=-2.75cm]current page.north west)
    {\sffamily\bfseries\fontsize{9}{11}\selectfont SERIE GUÍAS · TECNOLOGÍA E INFORMÁTICA};

  \begin{scope}[shift={([xshift=-3.05cm,yshift=-2.55cm]current page.north east)}]
    \fill[white,opacity=.80] (-.62,-.52) rectangle (.62,.52);
    \draw[milcGradeText,line width=1pt] (-.62,-.52) rectangle (.62,.52);
    \fill[milcGradeDark] (-.42,-.38) rectangle (-.22,.06);
    \fill[milcGradeAccent] (-.10,-.38) rectangle (.10,.24);
    \fill[milcGradePrimary!60!black] (.22,-.38) rectangle (.42,.38);
  \end{scope}

  \node[anchor=west,text=milcGradeText] at ([xshift=1.9cm,yshift=-12.85cm]current page.north west)
    {\sffamily\bfseries\fontsize{124}{124}\selectfont 7};
  % El circulito de grado (°) solo aplica a las guias de grado («11°»); en
  % Bebras y Semillero el numero grande es un momento/modulo, no un grado.
  \draw[milcGradeText,line width=1.7pt]
    ([xshift=4.52cm,yshift=-11.68cm]current page.north west) circle (.13cm);

  \node[anchor=west,text=milcGradeText,text width=8.7cm,align=left] at ([xshift=10.35cm,yshift=-11.6cm]current page.north west)
    {
     {\sffamily\bfseries\fontsize{19}{22}\selectfont Variables ---\\cajas con\\etiqueta\\que guardan\\información}\\[4mm]
     {\sffamily\bfseries\fontsize{23}{26}\selectfont Guía 15-7-TIC}\\[2mm]
     {\sffamily\fontsize{13}{16}\selectfont Período 2 · Algoritmia y pensamiento computacional}\\[5mm]
     {\sffamily\bfseries\fontsize{11}{13}\selectfont Institución Educativa Sor María Juliana}\\[1mm]
     {\sffamily\fontsize{10}{12}\selectfont Autor: Dr. Álvaro Cárdenas Orozco}
    };

  \path[fill=milcGradeSoft,rounded corners=7mm]
    ([xshift=1.35cm,yshift=.85cm]current page.south west) rectangle
    ([xshift=-1.35cm,yshift=4.25cm]current page.south east);
  \draw[milcGradeDark,line width=.65pt,rounded corners=7mm]
    ([xshift=1.35cm,yshift=.85cm]current page.south west) rectangle
    ([xshift=-1.35cm,yshift=4.25cm]current page.south east);
  \node[anchor=center,text=milcNegro,text width=17.0cm,align=center] at ([yshift=2.55cm]current page.south)
    {\sffamily\fontsize{10}{12}\selectfont Metodología MILC · Apertura ancestral · Pensamiento computacional · Triángulo Dussel-Estoicismo-Floridi\\
    \textcolor{milcGradeDark}{\bfseries Producto:} En el cuaderno: \textbf{(a) tabla de 10 variables} con \emph{nombre, tipo (número/texto/booleano), valor inicial} y \emph{ejemplo de uso}; \textbf{(b) un mini-algoritmo escrito} que use al menos 3 variables (ejemplo: \emph{``algoritmo del cumpleaños''} que use variables \texttt{edad}, \texttt{nombre}, \texttt{esCumpleHoy}); \textbf{(c) 5 reglas para nombrar variables} aplicadas. Firma.};
\end{tikzpicture}
\mbox{}
\newpage

\section*{Apertura: Variables --- cajas con etiqueta que guardan información}

\begin{softbox}{milcGradeDark}{milcGradeSoft}{Saber ancestral}
\textbf{Doña Sofía, la tendera del barrio de Cartago, tenía un sistema de cajitas en su tienda.} Cada cajita tenía una \textbf{etiqueta escrita a mano} con el nombre del producto: \emph{``arroz''}, \emph{``frijol''}, \emph{``café''}, \emph{``azúcar''}, \emph{``sal''}, \emph{``maíz molido''}. Adentro de cada cajita iba \emph{el producto correspondiente}. Cuando un cliente pedía \emph{``250 gramos de café''}, doña Sofía iba a la cajita \texttt{café}, sacaba 250 gramos, y entregaba. Si la cajita estaba vacía, doña Sofía la rellenaba. Si quería cambiar de un café por otro, vaciaba la cajita y ponía el nuevo. \textbf{Lo importante: la cajita era un lugar; lo que iba adentro podía cambiar}. Esa idea --- \emph{``un nombre fijo, un contenido variable''} --- es exactamente lo que en programación llamamos \textbf{variable}. Doña Sofía sabía gestionar variables sin haberlo estudiado. Hoy aprendes el concepto explícito.
\end{softbox}

\begin{softbox}{milcTurquesa}{milcGris}{Saber contemporáneo}
Una \textbf{variable} en programación es \emph{un nombre simbólico que guarda un valor que puede cambiar}. Imagínala como una \textbf{caja con etiqueta}: la etiqueta es el nombre (fijo), el contenido es el valor (variable). Cuando un algoritmo o programa necesita recordar un dato (la edad de alguien, el puntaje del juego, si la puerta está abierta), \textbf{crea una variable} y le pone un valor. Después puede \emph{leer} el valor (saber qué hay adentro), \emph{cambiar} el valor (vaciar y rellenar), o \emph{usar} el valor en operaciones. \textbf{Las variables son uno de los 3 componentes fundamentales} de cualquier programa (los otros 2 son: condicionales y bucles, que ves en las próximas 2 sesiones). \textbf{Hay 3 tipos básicos:} (1) \emph{número} (entero o decimal): \texttt{edad = 13}, \texttt{precio = 5.50}; (2) \emph{texto} (entre comillas): \texttt{nombre = ``María''}; (3) \emph{booleano} (solo dos valores: verdadero o falso): \texttt{tieneMascota = verdadero}.
\end{softbox}

\begin{softbox}{milcMostaza}{milcGris}{Pregunta-puente}
\textit{En tu vida diaria, ¿\textbf{qué información necesitas ``recordar''} cuando haces algo? Cuando juegas, hay un \emph{puntaje} (cambia). Cuando cocinas, hay un \emph{número de personas} (cambia). Cuando estudias, hay \emph{horas restantes} (cambia). Todas son variables.}
\end{softbox}

\begin{softbox}{milcVerde}{milcGris}{Saber y saber hacer}
Al terminar la clase: (1) podrás \textbf{identificar} qué es una variable y los 3 tipos básicos; (2) sabrás \textbf{aplicar} las 5 reglas para nombrar variables; (3) podrás \textbf{evaluar} si un algoritmo necesita variables o no; (4) habrás \textbf{creado} un mini-algoritmo que usa 3 variables.
\end{softbox}



\small
\begin{tabularx}{\textwidth}{>{\bfseries}p{3.0cm}Y}
\rowcolor{milcGradePrimary!12}Autor & Dr. Álvaro Cárdenas Orozco\\
DBA & Construye algoritmos y diagramas de flujo para resolver problemas paso a paso (MEN, T\&I 7°)\\
Referentes & Saber ancestral del tejedor · Pensamiento computacional · Scratch\\
Producto final & En el cuaderno: \textbf{(a) tabla de 10 variables} con \emph{nombre, tipo (número/texto/booleano), valor inicial} y \emph{ejemplo de uso}; \textbf{(b) un mini-algoritmo escrito} que use al menos 3 variables (ejemplo: \emph{``algoritmo del cumpleaños''} que use variables \texttt{edad}, \texttt{nombre}, \texttt{esCumpleHoy}); \textbf{(c) 5 reglas para nombrar variables} aplicadas. Firma.\\
\end{tabularx}
\normalsize

\puente{Hoy haces 4 cosas: aprendes qué es una variable, conoces los 3 tipos, aplicas las 5 reglas de nombrado, escribes un mini-algoritmo con variables.}

\begin{titlebox}{milcGradeDark}
Ruta MILC: así vamos a aprender
\end{titlebox}
\begin{tabularx}{\textwidth}{>{\centering\arraybackslash}X>{\centering\arraybackslash}X>{\centering\arraybackslash}X>{\centering\arraybackslash}X}
\phasecard{1}{milcVerde}{milcGris}{Escucha\\[-1mm]\small Observo, escucho y pregunto.} &
\phasecard{2}{milcTurquesa}{milcGris}{Sistematización\\[-1mm]\small Pensamiento computacional.} &
\phasecard{3}{milcMagenta}{milcGris}{Praxis\\[-1mm]\small Construyo y aplico.} &
\phasecard{4}{milcVino}{milcGris}{Evaluación\\[-1mm]\small Reflexión liberadora.}
\end{tabularx}


\puente{Empezamos con un juego: ¿qué información cambia en estas 6 situaciones?}

\begin{titlebox}{milcVerde}
Fase 1 · Escucha
\end{titlebox}

\begin{softbox}{milcVerde}{milcGris}{Escena para escuchar}
\textbf{Actividad 1 · IDENTIFICA --- ¿Qué información cambia?} (10 min · individual). Lee estas 6 situaciones y en el cuaderno \textbf{lista qué información cambia en cada una}. Cada información que cambie es una posible variable. \emph{(1) Estás jugando Roblox y miras el contador de monedas que aumenta.} \emph{(2) Tu mamá cocina una sopa para distinto número de personas cada vez.} \emph{(3) Estás haciendo una tarea y miras cuántas preguntas faltan.} \emph{(4) Es viernes y miras cuántos días faltan para tu cumpleaños.} \emph{(5) Llueve y revisas la temperatura cada hora.} \emph{(6) En clase llaman lista y se cuenta cuántos asistieron.}
\end{softbox}

\checkbox Leí las 6 situaciones.\\
\checkbox Identifiqué qué información cambia en cada una.\\
\checkbox Reconocí que esa información cambiante son variables.

\begin{infoband}{milcVerde}


\textbf{Lo que va al cuaderno (Actividad 1 · IDENTIFICA).} Título: «Actividad 1 --- Información que cambia». \textbf{Formato:} lista numerada del 1 al 6 con la variable identificada. \textbf{Extensión:} 6 líneas. \textbf{Sabes que terminaste cuando} reconoces que tu vida está llena de variables sin saberlo.
\end{infoband}

\puente{Después aprendes los 3 tipos + las 5 reglas para nombrar bien.}

\begin{titlebox}{milcTurquesa}
Fase 2 · Sistematización con pensamiento computacional
\end{titlebox}

\begin{softbox}{milcTurquesa}{milcGris}{Cuatro pilares aplicados al tema}
Lo que identificaste son \textbf{variables del mundo real}. En programación las nombramos explícitamente. Ahora aprendes los 3 tipos básicos + las 5 reglas para ponerles buen nombre.
\end{softbox}

\small
\begin{tabularx}{\textwidth}{>{\bfseries\RaggedRight\arraybackslash}p{3.6cm}Y}
\rowcolor{milcTurquesa!90}\color{white}Pilar & \color{white}\bfseries Aplicación al tema\\
1. Descomposición & \textbf{Tipo 1 --- Número.} Guarda un valor numérico. Puede ser \emph{entero} (sin decimales): \texttt{edad = 13}, \texttt{puntaje = 250}, \texttt{cantidadHermanos = 2}. O \emph{decimal} (con decimales): \texttt{altura = 1.65}, \texttt{precio = 5.50}, \texttt{notaPromedio = 4.2}. \textbf{Operaciones típicas:} sumar, restar, multiplicar, dividir, comparar (mayor/menor). \textbf{Tipo 2 --- Texto (string).} Guarda una cadena de caracteres. Va entre comillas: \texttt{nombre = ``María''}, \texttt{ciudad = ``Cartago''}, \texttt{materia = ``Tecnología''}. \textbf{Operaciones típicas:} unir, comparar, buscar dentro.\\
2. Reconocimiento de patrones & \textbf{Tipo 3 --- Booleano.} Guarda solo \textbf{dos valores}: \emph{verdadero} (true) o \emph{falso} (false). Sin matices. Ejemplo: \texttt{tieneMascota = verdadero}, \texttt{esCumpleHoy = falso}, \texttt{puertaAbierta = verdadero}. \textbf{Uso típico:} en condicionales (lo verás en S6). El nombre suele empezar con verbo \emph{es...}, \emph{tiene...}, \emph{está...} para que se entienda que es booleano. \textbf{Pista:} los booleanos son fundamentales en programación porque permiten tomar decisiones.\\
3. Abstracción & \textbf{Las 5 reglas para nombrar variables (universales en programación).} (1) \textbf{Empezar con minúscula:} \texttt{edad} sí, \texttt{Edad} mejor para clases (concepto avanzado). (2) \textbf{Camel case si son varias palabras:} \texttt{notaPromedio} (primera minúscula, siguientes con mayúscula). (3) \textbf{Nombre descriptivo, no críptico:} \texttt{x} no dice nada; \texttt{cantidadHermanos} sí dice. (4) \textbf{Sin espacios ni acentos:} \texttt{notaPromedio} sí; \texttt{nota promedio} no; \texttt{notaPromédio} mejor evítarlo. (5) \textbf{No empezar con número:} \texttt{2edad} no es válido; \texttt{edad2} sí.\\
4. Algoritmo conceptual & \textbf{Operaciones básicas con variables (en lenguaje natural).} \textbf{Asignar valor:} \texttt{edad = 13} (le pones 13 en la caja \texttt{edad}). \textbf{Leer valor:} cuando dices \texttt{edad}, te referís al 13 que está adentro. \textbf{Cambiar valor:} \texttt{edad = 14} (vacías la caja y le metes 14). \textbf{Usar en cálculo:} \texttt{edadProximaAno = edad + 1} (creas una nueva variable con el valor de \texttt{edad} más 1, o sea 14). En Scratch (S8) lo verás visualmente con bloques.\\
\end{tabularx}
\normalsize

\begin{drawbox}{milcTurquesa}{3 tipos + 5 reglas + operaciones básicas}
\textbf{3 tipos básicos:} \\[1mm]
\textbf{1.} NÚMERO (entero o decimal) --- 13, 1.65, 250. \\[1mm]
\textbf{2.} TEXTO (string) --- ``María'', ``Cartago''. \\[1mm]
\textbf{3.} BOOLEANO --- verdadero o falso solamente. \\[1mm]
\textbf{5 reglas:} minúscula inicial · camelCase · descriptivo · sin espacios/acentos · no empezar con número. \\[1mm]
\textbf{Operaciones:} asignar (=) · leer · cambiar · usar en cálculo.
\end{drawbox}

\begin{softbox}{milcMostaza}{milcGris}{Errores comunes y su corrección}
\textbf{Errores frecuentes al usar variables.} (1) \emph{Nombres incomprensibles:} \texttt{x}, \texttt{var1}, \texttt{tmp}. Cuando vuelvas en 2 días, no sabes qué significan. (2) \emph{Mezclar tipos:} sumar un texto con un número da error. \emph{``13''} (texto) más \texttt{2} (número) NO da 15 necesariamente. (3) \emph{Usar antes de asignar:} si pides leer \texttt{edad} sin haberle puesto valor, da error. (4) \emph{Olvidar las comillas en texto:} \texttt{nombre = María} sin comillas confunde al programa. (5) \emph{Crear demasiadas variables:} 50 variables se vuelven inmanejables. Una buena regla: solo crear variables que vas a usar al menos 2 veces.
\end{softbox}

\begin{infoband}{milcTurquesa}


\textbf{Lo que va al cuaderno (Actividad 2 · EXPLICA).} Título: «Actividad 2 --- 3 tipos + 5 reglas + operaciones». \textbf{Formato:} bloque A con tabla de los 3 tipos (tipo, ejemplo, operaciones típicas). Bloque B con las 5 reglas + 1 ejemplo malo y 1 bueno por regla. Bloque C con 4 operaciones básicas. \textbf{Extensión:} 12 líneas. \textbf{Sabes que terminaste cuando} podrías nombrar bien variables para cualquier algoritmo.
\end{infoband}

\puente{Con todo claro, escribes tu tabla de 10 variables + mini-algoritmo.}

\begin{titlebox}{milcMagenta}
Fase 3 · Praxis con pensamiento computacional
\end{titlebox}

\begin{softbox}{milcMagenta}{milcGris}{Construyendo el producto con los cuatro pilares}
\textbf{Actividad 3 · APLICA --- Tabla de 10 variables + mini-algoritmo con 3 variables} (28 min · individual). Vas a hacer 2 cosas. Primero, una tabla de 10 variables imaginadas. Segundo, escribir un mini-algoritmo que use 3 de ellas.
\end{softbox}

\small
\begin{tabularx}{\textwidth}{>{\bfseries\RaggedRight\arraybackslash}p{3.6cm}Y}
\rowcolor{milcMagenta!90}\color{white}Pilar & \color{white}\bfseries Aplicación al producto\\
Descomposición & \textbf{Paso 1 --- Tabla de 10 variables.} En una hoja del cuaderno, encabezado: \textbf{``Mis 10 variables''}. Tabla de 10 filas, 4 columnas: \emph{(1) Nombre de variable}, \emph{(2) Tipo (número/texto/booleano)}, \emph{(3) Valor inicial}, \emph{(4) Ejemplo de uso}. Llena las 10 filas con variables que podrías usar en programas reales. Sugerencias: \texttt{edadUsuario}, \texttt{nombreJugador}, \texttt{puntaje}, \texttt{vidasRestantes}, \texttt{tieneMascota}, \texttt{ciudadNatal}, \texttt{notaMatematicas}, \texttt{esFinDeSemana}, \texttt{horaActual}, \texttt{materiaPreferida}. Aplica las 5 reglas en cada nombre.\\
Patrones & \textbf{Paso 2 --- Escoge 3 variables para tu mini-algoritmo.} De las 10 que listaste, escoge 3 que se puedan combinar en un mini-algoritmo. Sugerencia útil: \texttt{nombre}, \texttt{edad}, \texttt{esCumpleHoy}.\\
Abstracción & \textbf{Paso 3 --- Escribe el algoritmo del cumpleaños.} Mini-algoritmo en lenguaje natural usando esas 3 variables: \emph{``ENTRADA: nombre (texto), edad (número), esCumpleHoy (booleano).''} \emph{``PASOS: 1. Mostrar saludo: 'Hola [nombre]'. 2. Mostrar: 'Tienes [edad] años'. 3. Si esCumpleHoy = verdadero, mostrar: 'Feliz cumpleaños, [nombre]! Ahora tienes [edad + 1] años'. Si esCumpleHoy = falso, mostrar: 'No es tu cumpleaños hoy'. 4. FIN.''}. Adáptalo a tu manera.\\
Algoritmo de armado & \textbf{Paso 4 --- Probar mentalmente el algoritmo.} Imagina que las variables tienen estos valores: \texttt{nombre = ``Lara''}, \texttt{edad = 12}, \texttt{esCumpleHoy = verdadero}. \emph{Ejecuta} el algoritmo paso a paso en tu cabeza y escribe qué se mostraría. Deberías obtener: \emph{``Hola Lara. Tienes 12 años. Feliz cumpleaños, Lara! Ahora tienes 13 años.''}. Si tu salida es distinta, revisa el algoritmo.\\
\end{tabularx}
\normalsize

\begin{drawbox}{milcMagenta}{Antes de cerrar la S5}
\checkbox Tabla de 10 variables con 4 columnas. \\[1mm]
\checkbox Cada nombre sigue las 5 reglas (minúscula, camelCase, etc.). \\[1mm]
\checkbox Cada variable tiene tipo correctamente identificado. \\[1mm]
\checkbox Mini-algoritmo usa 3 variables. \\[1mm]
\checkbox Probé mentalmente el algoritmo con valores concretos. \\[1mm]
\checkbox Reflexión y firma están al final.
\end{drawbox}

\begin{softbox}{milcMagenta}{milcGris}{Plantilla de trabajo}
\textbf{Estructura.} \textit{Paso 1:} tabla 10 × 4 (nombre, tipo, valor inicial, ejemplo). \textit{Paso 2:} escoger 3 variables. \textit{Paso 3:} algoritmo en lenguaje natural usando las 3. \textit{Paso 4:} ejecución mental con valores concretos. \textit{Paso 5:} verificar 5 reglas. \textit{Paso 6:} reflexión + firma.
\end{softbox}

\begin{infoband}{milcMagenta}


\textbf{Lo que va al cuaderno (Actividad 3 · APLICA).} Título: «Actividad 3 --- 10 variables + mini-algoritmo». \textbf{Formato:} tabla 10x4 + algoritmo + prueba mental + reflexión. \textbf{Extensión:} 1 página completa. \textbf{Sabes que terminaste cuando} podrías mostrar al profe tus variables y él entendería para qué sirven solo leyendo los nombres.
\end{infoband}

\puente{El producto es la tabla + algoritmo + 5 reglas firmadas.}

\begin{titlebox}{milcMostaza}
Producto de la sesión
\end{titlebox}
\textbf{10 variables nombradas + mini-algoritmo con 3 variables · 5 reglas aplicadas}.
\begin{softbox}{milcMostaza}{milcGris}{Criterios de calidad}
(1) Tabla con 10 variables (nombre, tipo, valor inicial, ejemplo). (2) Las 5 reglas de nombrado están aplicadas correctamente. (3) El mini-algoritmo usa al menos 3 variables. (4) Se probó mentalmente el algoritmo con valores concretos. (5) Hay reflexión final firmada.
\end{softbox}

\puente{Antes de salir, intercambia con un compañero: ¿tus nombres de variable son claros?}

\begin{titlebox}{milcVino}
Fase 4 · Evaluación liberadora — cinco dimensiones
\end{titlebox}

\begin{softbox}{milcVino}{milcGris}{Sentido de la evaluación}
Evaluar no es solo revisar una nota. Es reconocer qué aprendí, qué cambió en mí, qué emociones manejé, cómo me relaciono con la ciudad y la comunidad, y qué le dejaré a quien viene detrás.
\end{softbox}

\textbf{Mi reflexión liberadora en cinco dimensiones.} Completa cada frase iniciada:

\small
\begin{tabularx}{\textwidth}{>{\bfseries\RaggedRight\arraybackslash}p{3.4cm}Y>{\RaggedRight\arraybackslash}p{4.6cm}}
\rowcolor{milcVino}\color{white}Dimensión & \color{white}\bfseries Mi respuesta (frame starter) & \color{white}\bfseries Algo que descubrí\\
1. Personal & Lo que más cambió en mí fue \linead{6.5cm} & \linead{4.2cm}\\
2. Emocional & La emoción más fuerte fue \linead{2.5cm} y la regulé \linead{3cm} & \linead{4.2cm}\\
3. Ciudadana & En democracia esto importa porque \linead{6cm} & \linead{4.2cm}\\
4. Local (Cartago) & En mi barrio sería útil para \linead{6.5cm} & \linead{4.2cm}\\
5. Intergeneracional & A mi abuela/abuelo le diría \linead{6.5cm} & \linead{4.2cm}\\
\end{tabularx}
\normalsize

\begin{infoband}{milcVino}
\textbf{Para la próxima sustentación.} Vuelve a esta tabla en seis meses. Verás que las respuestas cambian — no porque lo aprendido cambie, sino porque tú cambias. La evaluación liberadora es una conversación contigo a través del tiempo.
\end{infoband}


\puente{Cerramos con 3 voces sobre el poder de nombrar las cosas.}

\begin{titlebox}{milcGradeDark}
Triángulo de pensamiento · cierre filosófico
\end{titlebox}

\begin{softbox}{milcVino}{milcGris}{Dussel · lente del nosotros}
\textit{````Nombrar las cosas correctamente es el primer acto de pensamiento claro.''''}

Las variables se llaman \emph{``variables''} porque varían (cambian). Pero su \textbf{nombre es fijo}: \emph{ese nombre es el primer acto de pensamiento}. Si nombras una variable \texttt{x}, no piensas claro. Si la nombras \texttt{cantidadMonedasGanadas}, piensas claro. \emph{La calidad del nombre refleja la calidad del pensamiento}. Doña Sofía con sus cajitas etiquetadas \emph{café}, \emph{arroz}, \emph{azúcar} también lo entendía. Las cajas sin etiqueta no sirven; los nombres dan función.

\textbf{¿Cuántas cosas en mi vida tengo sin nombrar bien y que por eso me confunden?}
\end{softbox}
\linead{\linewidth}\\[1mm]
\linead{\linewidth}

\begin{softbox}{milcMostaza}{milcGris}{Estoicismo · Marco Aurelio · cuidado interior}
\textit{````La palabra exacta resuelve confusiones que las explicaciones largas no logran.''''}

En programación, \textbf{el nombre de una variable es comunicación con tu yo del futuro}. Vuelves a leer tu código en 6 meses, no recuerdas detalles, pero leyendo \texttt{cantidadMonedasGanadas} entiendes inmediatamente qué guarda. Si lo hubieras llamado \texttt{x}, tendrías que rastrear el código para entender. \emph{Las 5 reglas de nombrado son disciplina lingüística antigua trasladada a programación}.

\textbf{¿Cuando escribo (texto, código, listas), mi yo del futuro lee con claridad o tiene que adivinar?}
\end{softbox}
\linead{\linewidth}\\[1mm]
\linead{\linewidth}

\begin{softbox}{milcTurquesa}{milcGris}{Floridi · lente de la infoesfera}
\textit{````El 50\% del oficio de programar es nombrar bien. El otro 50\% es lógica. Los buenos programadores invierten en ambos.''''}

Hay un dicho en programación: \emph{``solo hay dos cosas difíciles: nombrar bien y manejar la caché''}. Es bromas, pero verdadero. \textbf{Aprender hoy a nombrar variables con criterio te coloca en el camino correcto del oficio}. Quien nombra mal escribe código que nadie (ni el mismo autor) entiende. Quien nombra bien produce código que dura años legible.

\textbf{¿Voy a ser un programador (o pensador) que nombra bien, o uno que mete \texttt{x} en todas partes?}
\end{softbox}
\linead{\linewidth}\\[1mm]
\linead{\linewidth}

\begin{titlebox}{milcVino}
Compromiso final
\end{titlebox}

\small
\begin{tabularx}{\textwidth}{>{\RaggedRight\arraybackslash}p{3.0cm}YYY}
\rowcolor{milcVino}\color{white}\bfseries Criterio & \color{white}\bfseries Lo logré & \color{white}\bfseries En proceso & \color{white}\bfseries Necesito apoyo\\
Comprensión del tema & Explico ideas centrales con mis palabras. & Reconozco ideas pero falta precisión. & Necesito repasar conceptos base.\\
Pensamiento computacional & Apliqué los 4 pilares al diseñar mi producto. & Apliqué algunos pilares. & Necesito releer la fase 2.\\
Aplicación y producto & Mi producto cumple los criterios y se puede revisar. & Producto funcional con ajustes pendientes. & Aún en construcción.\\
Reflexión 5 dimensiones & Respondí cada dimensión con honestidad. & Respondí algunas con detalle. & Mis respuestas son muy generales.\\
Triángulo de pensamiento & Conecté las 3 voces con mi trabajo real. & Conecté al menos una. & No relaciono las voces con mi proceso.\\
\end{tabularx}
\normalsize

\textbf{Hoy entendí que:}\\[1mm]
\linead{\linewidth}\\[1mm]
\linead{\linewidth}

\textbf{Mi compromiso para la próxima sesión es:}\\[1mm]
\linead{\linewidth}\\[1mm]
\linead{\linewidth}

\begin{softbox}{milcMostaza}{milcGris}{Cita de despedida · Marco Aurelio}
\textit{``El obstáculo en el camino se vuelve el camino. Lo que estorba al camino se vuelve camino.''}\\[1mm]
Cada error de hoy, bien observado, se convierte en pieza del camino que pisarás mañana.
\end{softbox}

\small
\begin{tabularx}{\textwidth}{>{\bfseries}p{3.6cm}Y}
\rowcolor{milcGradeDark!12}Nombre completo: & \linead{10cm}\\
Fecha: & \linead{4cm} \quad Curso/grupo: \linead{4cm}\\
Mi firma: & \linead{9cm}\\
\end{tabularx}
\normalsize

\begin{infoband}{milcGradeDark}
\textbf{Para la próxima sesión.} Esta semana, identifica variables en al menos 3 situaciones de tu vida real (un juego, una receta, una rutina). Anota nombres siguiendo las 5 reglas. La próxima clase aprendes \textbf{condicionales (si... entonces)}: las decisiones que usan los algoritmos para escoger caminos según las variables.
\end{infoband}

\end{document}
